% !TEX root = statpointmenger.tex
\section{\secA}

Before we begin the discussion of the first variation, let us rewrite the integral Menger curvature using
the symmetry and a suitable covering of the domain of integration $(\R/\Z)^3$ by domains, on which it is 
easier to estimate the terms that will appear.
The general idea here is quite similar
to~\cite{blatt:imcc} and Hermes~\cite{hermes}, but we will 
show that it is actually enough to integrate over a certain subdomain of $(\mathbb R / \mathbb Z)^3$.

To this end we define the range of integration
\begin{equation}\label{eq:deco}
 D:=\sets{(v,w) \in (-\tfrac12,0) \times (0,\tfrac12)}{w\leq 1+2v, v \geq -1 + 2w},
\end{equation}
which is depicted in Figure~\ref{fig:domain}.

\begin{figure}[h]
 \centering
 \includegraphics{fig-2.pdf}
  \caption{The range of integration $D$}
 \label{fig:domain}
\end{figure}

\begin{lemma}[Domain decomposition]\label{lem:domdec}
 Let $f\in L^{1}\br{\br{\R/\Z}^{3}}$ be symmetric in all components,
 i.e.\@ $f=f\circ\s$ for all permutations $\s\in\mathfrak S_{3}$.
 Then
 \[ \iiint_{(\R/\Z)^{3}} f(u_{1},u_{2},u_{3})\d u_{1}\d u_{2}\d u_{3} = 6\iiint_{\R/\Z\times D} f(u,u+v,u+w)\d w\d v\d u. \]
\end{lemma}

\begin{proof}
 Let $P_{\s}\in\R^{3\times 3}$ denote the permutation matrix corresponding to $\s\in\mathfrak S_{3}$. 
 We first show that the images
 $\sets{P_{\s}\br{\R/\Z\times D}}{\s\in\mathfrak S_{3}}$
 cover $(\R/\Z)^{3}$.

 Consider $(u_{1},u_{2},u_{3})\in\br{\R/\Z}^{3}$. Then after a suitable permutation we can assume that
 \begin{equation*}
  d_{\mathbb R / \mathbb Z}(u_1,u_3) = \max\br{d_{\mathbb R / \mathbb Z}(u_1,u_2),d_{\mathbb R / \mathbb Z}(u_2,u_3),
  d_{\mathbb R / \mathbb Z}(u_1, u_3) }
 \end{equation*}
 where $d_{\mathbb R / \mathbb Z}(x,y) = \min\{|x-y-k|: k \in \mathbb Z\}
 \in[0,\frac12]$ denotes the
 distance in $\mathbb R / \mathbb Z$.
 Hence, interchanging $u_{1}$ and $u_{3}$ if necessary,
 there are $(v,w) \in \bv{-\frac12,0}\times\bv{0,\frac12}$ with
 $$
   u_1 = u_2 + v,\qquad u_3 = u_2 + w,
 $$
 and 
 \begin{align*} \max\br{-v,w} &=\max\br{d_{\R/\Z}(u_{1},u_{2}),d_{\R/\Z}(u_{2},u_{3})}
    \\ &\le d_{\R/\Z}(u_{1},u_{3}) = \min\br{w-v,1-(w-v)}\le1-w+v, \end{align*}
%  d_{\mathbb R / \mathbb Z}(u_1,u_3) = 1- (w-v) \leq \max\{-v,w\} = \max\{d_{\mathbb R / \mathbb Z}(u_1,u_2),d_{\mathbb R / \mathbb Z}(u_2,u_3),
%  d_{\mathbb R / \mathbb Z}(u_1, u_3) \},
 %But this implies that
 so $(v,w) \in\overline D$.
 Since furthermore 
 \begin{equation*}
\#\mathfrak S_3 \cdot \abs{\R/\Z\times D}=6 \cdot \tfrac16 = \abs{(\mathbb R / \mathbb Z)^3}
 \end{equation*}
 the sets $\sets{P_{\s}\br{\R/\Z\times D}}{\s\in\mathfrak S_{3}}$
 form up to sets of measure zero a disjoint partition of $(\R/\Z)^{3}$.
\end{proof}


Following Lemma~\ref{lem:domdec} we derive using~\eqref{eq:short-notation}
\begin{equation}\label{eq:intM-sn''}
 \E(\g)=6\iiint_{\R/\Z\times D}
 \frac{\abs{\T v0\wedge\T 0w}^{q}}{\br{\abs{\T v0}\abs{\T 0w}\abs{\T vw}}^p}
 {\abs{\dg(u)}\abs{\dg(u+v)}\abs{\dg(u+w)}}
 {\d w\d v\d u}.
\end{equation}

\begin{proof}[Theorem~\ref{thm:energy-space}]
Recall that any embedded $\Wia[1+s,\rho]$-curve, $s>\tfrac1\rho$, is bi-Lipschitz continuous~\cite[Lemma~2.1]{blatt:bre}.
We obtain, using $\abs{a\wedge b}=\abs{a\wedge(a\pm b)}\le\abs a\abs{a\pm b}$ for $a,b\in\R^{3}$,
\begin{align*}
 \E(\g) &\le C \iiint_{\R/\Z\times D}\frac{\abs{\int_0^1\dg(u+\theta_{1} v)\d\theta_{1}\wedge
 \int_0^1\dg(u+\theta_{2} w)\d\theta_{2}}^{q}}{\abs v^{p-q}\abs w^{p-q}\abs{v-w}^{p}} \d w\d v\d u\\
 &\le C \tint\frac{\int_0^1\abs{\dg(u+\theta v)-\dg(u+\theta w)}^{q}\d\theta}{\abs v^{p-q}\abs w^{p-q}\abs{v-w}^{p}} \d w\d v\d u\\
 &\le C \int_0^1\tint\frac{\abs{\dg(u+\theta (v-w))-\dg(u)}^{q}}{\abs v^{p-q}\abs w^{p-q}\abs{v-w}^{p}} \d w\d v\d u\d\theta
\end{align*}
where $C$ depends on $p,q$ and $\g$.
Substituting
$\Phi:(v,w)\mapsto (t,\tilde w) := \br{\frac v{v-w},\theta(v-w)}$,
$\abs{\det D\Phi(v,w)} = \frac\theta{\abs{v-w}}$,
$\Phi(D)\subset[0,1]\times[-1,0]$,
we arrive at
\begin{equation}\label{eq:intM<seminorm}
\begin{split}
 &\E(\g) \\
 &\le C \int_{0}^{1}\theta^{-1}\int_{\R/\Z}\int_{0}^{1}\int_{-1}^{0}
 \frac{\abs{\dg(u+\tilde w)-\dg(u)}^{q}}{\abs {t\frac{\tilde w}\theta}^{p-q}\abs{(t-1)\frac{\tilde w}\theta}^{p-q}\abs{\frac{\tilde w}\theta}^{p-1}} \d\tilde w\d t\d u\d\theta \\
 &\le C \int_{0}^{1}{\theta^{3p-2q-2}}\d\theta
 \int_{0}^{1}\frac{\d t}{\abs{t(1-t)}^{p-q}}
 \int_{\R/\Z}\int_{-1}^{0}
 \frac{\abs{\dg(u+\tilde w)-\dg(u)}^{q}}{\abs {\tilde w}^{3p-2q-1}} \d\tilde w\d u \\
 &\le C\br{\seminorm{\g}_{\W}^{q}+\lnorm[\infty]\dg^{q}}
 \le C\norm{\g}_{\W}^{q}.
\end{split}
\end{equation}
For the other implication, we first derive for given vectors
$a,b\in\R^{n}$, $\abs a=\abs b=1$, $\sp{a,b}\ge0$,
\begin{equation}\label{eq:a-wedge-b}
 \abs{a\wedge b}^{2} = \abs a^{2}\abs b^{2}-\sp{a,b}^{2}
   \ge1-\sp{a,b}=\tfrac12\abs{a-b}^{2}.
\end{equation}
By uniform continuity of $\dg$
we may choose $\delta=\delta(\g)\in(0,\tfrac12)$ such that
\begin{equation}
 \abs{\dg(u+v)-\dg(u+w)} \le \tfrac1{10}
 \qquad \text{for all } u\in\R/\Z, \quad v,w\in[-\delta,\delta].
\end{equation}
In fact, we may choose $\delta$ to be maximal, i.e.\@ we
assume that there are $\tilde u\in\R/\Z$, $\tilde v,\tilde w\in[-\delta,\delta]$
with
\begin{equation}\label{eq:delta=max}
 \abs{\dg(\tilde u+\tilde v)-\dg(\tilde u+\tilde w)} = \tfrac1{10}.
\end{equation}
We fix $u_{0}\in\R/\Z$. As $\g\in C^{1}\rzd$ we may apply a suitable
translation and rotation of the ambient space~$\R^{n}$ such that
$\g(u_{0})=0$ and there is a function $f\in C^{1}(\R,\R^{n-1})$
with $\lnorm[\infty]f\le1$ and $f(0)=0$ such that
$\tilde\g(u):=(u,f(u))$ satisfies $\tilde\g(B_{2\delta}(0))\subset \g(\R/\Z)$.
Then
\begin{equation}\label{eq:tilde-bilip}
 \tfrac12\abs{\T[\tilde\g]v0}\le\abs{v}\le\abs{\T[\tilde\g]v0}.
\end{equation}

Arc-length parametrization of $\g$ gives
\renewcommand{\T}[3][\tg]{\triangle_{#2,#3}{#1}}
\begin{equation}\label{eq:tilde-g}
\begin{split}
 &\E(\tilde\g) \\
 &\ge c\int_{u_0-\delta}^{u_0+\delta}\int_{-\delta}^\delta\int_{-\delta}^\delta
 \frac{\abs{\fracabs{\T v0}\wedge\fracabs{\T 0w}}^{q}}{{\abs{\T v0}^{p-q}\abs{\T 0w}^{p-q}\abs{\T vw}^p}}
 {\d w\d v\d u} \\
 &\refeq[\ge]{a-wedge-b} c\int_{u_0-\delta}^{u_0+\delta}\int_{-\delta}^\delta\int_{-\delta}^\delta
 \frac{\abs{\sign v\fracabs{\T v0}+\sign w\fracabs{\T 0w}}^{q}}{{\abs{\T v0}^{p-q}\abs{\T 0w}^{p-q}\abs{\T vw}^p}}
 {\d w\d v\d u} \\
 &\refeq[\ge]{tilde-bilip} c\int_{u_0-\delta}^{u_0+\delta}\int_{-\delta}^\delta\int_{-\delta}^\delta
 \frac{\abs{\sign v\fracabs{\T v0}+\sign w\fracabs{\T 0w}}^{q}}{{\abs{v}^{p-q}\abs{w}^{p-q}\abs{v-w}^p}}
 {\d w\d v\d u} \\
 &\ge c\int_{u_0-\delta}^{u_0+\delta}\int_{-\delta}^\delta\int_{-\delta}^\delta
 \br{\frac{\abs{\sign v\fracabs{\T v0}+\sign w\fracabs{\T 0w}}^{q}}{{\abs{v}^{p-q}\abs{w}^{p-q}\abs{v-w}^p}}
 + \frac{\abs{-\sign v\fracabs{\T {-v}0}+\sign w\fracabs{\T 0w}}^{q}}{{\abs{v}^{p-q}\abs{w}^{p-q}\abs{v+w}^p}}}
 {\d w\d v\d u} \\
 &\ge c\int_{u_0-\delta}^{u_0+\delta}\int_{-\delta}^{\delta}\int_{-\abs v}^{\abs v}
 \br{\frac{\abs{\sign v\fracabs{\T v0}+\sign w\fracabs{\T 0w}}^{q}}{{\abs{v}^{p-q}\abs{w}^{p-q}\abs{v+w}^p}}
 +\frac{\abs{-\sign v\fracabs{\T {-v}0}+\sign w\fracabs{\T 0w}}^{q}}{{\abs{v}^{p-q}\abs{w}^{p-q}\abs{v-w}^p}}}
 {\d w\d v\d u} \\
 &\ge c\int_{u_0-\delta}^{u_0+\delta}\int_{-\delta}^{\delta}\int_{-\abs v}^{\abs v}
 \frac{\abs{\sign v\fracabs{\T v0}+\sign w\fracabs{\T 0w}}^{q}
 +\abs{-\sign v\fracabs{\T {-v}0}+\sign w\fracabs{\T 0w}}^{q}}{{\abs{v}^{3p-2q}}}
 {\d w\d v\d u} \\
 &\ge c\int_{u_0-\delta}^{u_0+\delta}\int_{-\delta}^{\delta}
 \frac{\abs{\fracabs{\T v0}+{\fracabs{\T {-v}0}}}^{q}}{{\abs{v}^{3p-2q-1}}}
 {\d v\d u}
% &\ge c\int_{u_0-\delta}^{u_0+\delta}\int_{-\delta}^{\delta}
% \frac{\abs{{\T v0}-{{\T {-v}0}}}^{q}}{{\abs{v}^{3p-q-1}}}
% {\d v\d u}
% -C\int_{u_0-\delta}^{u_0+\delta}\int_{-\delta}^{\delta}
% \frac{\abs{\frac1{\abs{\T v0}}-\frac1{\abs{\T {-v}0}}}^{q}}{{\abs{v}^{3p-2q-1}}}
% {\d v\d u} \\
% 
% 
% &\ge c\int_{u_0-\delta}^{u_0+\delta}\int_{-delta}^0\int_0^\delta
% \frac{\abs{\fracabs{\T v0}+\fracabs{\T 0w}}^{q}+\abs{\fracabs{\T v0}+\fracabs{\T 0w}}^{q}}{{\abs{v}^{p-q}\abs{w}^{p-q}\abs{v+w}^p}}
% {\d w\d v\d u} \\
% &\ge c\int_{u_0-\delta}^{u_0+\delta}\int_{0}^{\delta}\int_{0}^{v/2}
% \frac{\abs{\frac{\T vw}{\abs{\T v0}}+\T 0w\br{\frac1{\abs{\T v0}}-\frac1{\abs{\T 0w}}}}^{q}}{{{v}^{p-q}{w}^{p-q}\br{v+w}^p}}
% {\d w\d v\d u} \\
\end{split}
\end{equation}
where $c>0$ only depends on $p$ and $q$.
The last line in~\eqref{eq:tilde-g} is bounded below by
\begin{align*}
 &c\int_{u_0-\delta}^{u_0+\delta}\int_{-\delta}^{\delta}
 \frac{\abs{{\T v0}+{{\T {-v}0}}}^{q}}{{\abs{v}^{3p-q-1}}}
 {\d v\d u}
 -C\int_{u_0-\delta}^{u_0+\delta}\int_{-\delta}^{\delta}
 \frac{\abs{{\T{-v}0}}^{q}\abs{\frac1{\abs{\T v0}}-\frac1{\abs{\T {-v}0}}}^{q}}{{\abs{v}^{3p-2q-1}}}
 {\d v\d u} \\
 &\refeq[\ge]{tilde-bilip} c{
 \int_{u_{0}-\delta}^{u_{0}+\delta}\int_{-\delta}^{\delta}
 \frac{\abs{\tg(u+v)-2\tg(u)+\tg(u-v)}^{q}}{\abs v^{3p-q-1}}\d v\d u
 }
 -C\int_{u_0-\delta}^{u_0+\delta}\int_{-\delta}^{\delta}
 \frac{\abs{\frac v{\abs{\T v0}}-\frac{v}{\abs{\T {-v}0}}}^{q}}{{\abs{v}^{3p-2q-1}}}
 {\d v\d u}.
\end{align*}
By
\[ \abs{\frac v{\abs{\T v0}}-\frac{v}{\abs{\T {-v}0}}}
   \le \abs{\frac{(v,\T[f]v0)}{\abs{\T v0}}+\frac{(-v,\T[f]{-v}0)}{\abs{\T {-v}0}}}
   = \abs{\fracabs{\T v0}+\fracabs{\T {-v}0}} \]
we may use~\eqref{eq:tilde-g} to absorb the last term
which finally leads to
\[ \E(\g)\ge c{
 \int_{u_{0}-\delta}^{u_{0}+\delta}\int_{-\delta}^{\delta}
 \frac{\abs{\tg(u+v)-2\tg(u)+\tg(u-v)}^{q}}{\abs v^{3p-q-1}}\d v\d u
 }. \]
%\renewcommand{\T}[3][\g]{\tilde\triangle_{#2,#3}{#1}}%
%
Since reparametrization to arc-length preserves regularity, we arrive at
\begin{equation}\label{eq:delta-est}
 \E(\g)\ge c{
 \int_{u_{0}-\delta}^{u_{0}+\delta}\int_{-\delta}^{\delta}
 \frac{\abs{\g(u+v)-2\g(u)+\g(u-v)}^{q}}{\abs v^{3p-q-1}}\d v\d u
 }.
\end{equation}
As $u_{0}$ was chosen arbitrarily, we obtain
\begin{equation}\label{eq:semi-intM}
 {\seminormv{\g}}_{\W[(3p-2)/q-1,q]}^{q}
 \le C\br{\E(\g)+\lnorm[\infty]{\dg}^{q}\delta^{-3p+2q+2}}
\end{equation}
uniformly on $\R/\Z$.
Since the exponent $-3p+2q+2$ is negative,
we have to show that $\delta$ is uniformly
bounded away from zero in order to finish the proof.
To this end we will establish the Morrey-type estimate
 \begin{equation}\label{eq:Campanato}
  \lnorm[\infty]{\gamma'(\cdot+w) - \gamma'(\cdot)} \leq C \E(\g)^{1/q} |w|^{\alpha}
  \qquad\text{for all }w\in[-\tfrac12,\tfrac12]
 \end{equation}
 where $\alpha = 3(p-1)/q-2>0$.
 As $\delta$ was chosen to be maximal with respect to~\eqref{eq:delta=max},
 we arrive at
 \begin{equation*}
  \tfrac1{10} \le C \E(\g)^{1/q} \delta^{\alpha}
 \end{equation*}
 which, applied to~\eqref{eq:semi-intM}, gives~\eqref{eq:energy-space}.

It remains to prove~\eqref{eq:Campanato} which follows
by standard arguments due to Campanato~\cite{campanato}.
Let $\g'_{B_r (x)}$ denote the integral mean 
 of $\g'$ over $B_r(x)$. We calculate for $x \in \mathbb R / \mathbb Z$
 and $r \in (0,\delta)$ 
 \begin{align*}
  \frac 1 {2r}\int_{B_r(x)} |\gamma'(v)- \g'_{B_r(x)}| \d v
  &\leq \frac 1 {4r^2} \int_{B_r(x)} \int_{B_r(x)}| \gamma'(v) - \gamma'(u)| \d u \d v
  \\
  &\leq  \left(\frac 1 {4r^2} \int_{B_r(x)} \int_{B_r(x)}| \gamma'(v) - \gamma'(u)|^q \d u \d v\right)^{1/q}
  \\
  &\leq C r^{\alpha} \left(\int_{B_r(x)} \int_{B_r(x)} \frac{| \gamma'(v) - \gamma'(u)|^q} {|u-v|^{3p-2q-1}}\d u \d v\right)^{1/q} 
  \\
  &\refeq[\le]{delta-est} C r^{\alpha} \E(\g)^{1/q}.
 \end{align*}
 As~\eqref{eq:Campanato} only involves the domain of $\g$
 up to a measure zero set, we may restrict to Lebesgue points.
 We choose two Lebesgue points $u,v \in \mathbb R / \Z$ of $\gamma'$
 with $r:=|u-v| \in (0,\tfrac\delta2)$. Then
 \begin{equation*}
  |\dg(u) - \dg(v)| 
  \leq \sum_{k=0}^\infty \left| \dg_{B_{2^{1-k}r}(u)} - \dg_{B_{2^{-k}r}(u)} \right|
  + \left| \dg_{B_{2r}(u)} - \dg_{B_{2r}(v)} \right|
  + \sum_{k=0}^\infty \left| \dg_{B_{2^{1-k}r}(v)} - \dg_{B_{2^{-k}r}(v)} \right|.
  \end{equation*}
  Since
  \begin{align*}
   \left| \dg_{B_{2r}(u)} - \dg_{B_{2r}(v)} \right| 
   &\leq \frac {\int_{B_{2r}(u)} |\dg(x) - \dg_{B_{2r}(u)}| \d x +
    \int_{B_{2r}(v)} |\dg(x) - \dg_{B_{2r}(v)}| \d x} {|B_{2r} (u) \cap B_{2r}(v)|}
    \\
    &\leq C |u-v|^\alpha \E(\g)^{1/q} 
  \end{align*}
  as $r=|u-v|$ and, for all $y \in \mathbb R / \mathbb Z$, $R\in(0,\tfrac\delta2)$,
  \begin{align*}
   \left| \dg_{B_{2R}(y)} - \dg_{B_{R}(y)} \right| 
   &\leq \frac {\int_{B_{R}(y)} |\dg(x) - \dg_{B_{2R}(y)}| \d x +
    \int_{B_{R}(y)} |\dg(x) - \dg_{B_{R}(y)}| \d x} {2R}
    \\
    &\leq C R^\alpha \E(\g)^{1/q},
  \end{align*}
  we deduce
%  \begin{align*}
   $|\dg(u) - \dg(v)| 
   \leq C \left( \sum_{k=0} ^\infty 2^{-k\alpha} + 1 +  \sum_{k=0} ^\infty 2^{-k\alpha}   \right)
   |u-v|^\alpha \E(\g)^{1/q}$.
%  \end{align*}
  Thus
%  \begin{equation*}
   $|\dg(u) - \dg(v)| \leq C |u-v|^\alpha \E(\g)^{1/q}$
%  \end{equation*}
 for all Lebesgue points of $\gamma'$ with $|u-v| < \frac\delta2$.
% Since Lebesgue points are dense and using the triangle inequality this 
% proves \eqref{eq:Campanato}.
 The case $|u-v| \geq \frac\delta2$ follows by the triangle inequality.
\end{proof}

Let us conclude this section by briefly commenting on the other
ranges in the $(p,q)$-domain, see Figure~\ref{fig:range}.

\begin{remark}[Non-repulsive energies for $p<\frac23q+1$]\label{rem:non-rep}
 A bi-Lipschitz estimate is not guaranteed for injective curves if $p<\frac23q+1$. We briefly give the following example.
 Consider the curves $u\mapsto (u,0,0)$ and $u\mapsto (0,u,\delta)$ for $u\in[-1,1]$, $\delta\in[0,1]$.
 The interaction of these strands leads to the $\E$-value
 \begin{align*}
  &C\iiint_{[-1,1]^3} \frac{\br{\delta^{2}+u^{2}}^{q/2}}{\abs{v-w}^{p-q}{{\br{\delta^{2}+u^2+v^{2}}^{p/2}\br{\delta^{2}+u^2+w^{2}}^{p/2}}}}\d w\d v\d u \\
  &\le C\iiint_{[-1,1]^3} \frac{\br{\delta^{2}+u^{2}}^{(q-p)/2}}{\abs{v-w}^{p-q}{{\br{\delta^{2}+u^2+v^{2}+w^{2}}^{p/2}}}}\d w\d v\d u.
 \end{align*}
 Introducing polar coordinates $u=r\cos\th$, $v=r\sin\th\cos\varphi$,
 $w=r\sin\th\sin\varphi$, the former quantity is bounded by
 \begin{align*}
  & C\int_0^{\sqrt3}\int_0^{\pi}
  \frac{\br{\delta^{2}+r^{2}\cos^{2}\th}^{(q-p)/2}r^2\sin\th}
  {r^{p-q}\sin^{p-q}\th\br{\delta^{2}+r^{2}}^{p/2}}
  \d\th\d r
  \underbrace{\int_{0}^{2\pi}\frac{\d\varphi}{\abs{\cos\varphi-\sin\varphi}^{p-q}}}_{{\le C}} \\
  & \le C\int\limits_0^{\sqrt3}\br{\int\limits_{[0,\frac\pi4]\cup[\frac{3\pi}4,\pi]}
  \frac{\br{\delta^{2}+r^{2}}^{(q-p)/2}\d\th}
  {r^{p-q-2}\underbrace{\sin^{p-q-1}\th}_{\ge1}\br{\delta^{2}+r^{2}}^{p/2}}
  +\int\limits_{\frac\pi4}^{\frac{3\pi}4}
  \frac{\br{\delta^{2}+r^{2}\cos^{2}\th}^{(q-p)/2}r\sin\th}
  {r^{p-q-1}\sin^{p-q}\th\br{\delta^{2}+r^{2}}^{p/2}}
  \d\th
  }\d r \\
  & \le C\int_0^{\sqrt3}\br{
  \br{\delta+r}^{-3p+2q+2}
  +r^{-p+q+1}\br{\delta+r}^{-p}
  \int_{0}^{r}\br{\delta^{2}+\sigma^{2}}^{(q-p)/2}\d\sigma
   }\d r \\
  &\le C\br{1-\delta^{-3p+2q+3}}\le C.
% \label{eq:non-rep}
 \end{align*}
 Using Theorem~\ref{thm:energy-space}
 and the monotonicity of $\E[\cdot,q]$ for fixed $q$, it is easy to
 produce a family of knots uniformly converging to a non-embedded curve
 without an energy blow-up as $\delta\searrow0$,
 so these energies are not self-repulsive.
\end{remark}

\begin{remark}[Singular energies for $p\ge q+\tfrac23$, $q>1$]\label{rem:singular}
 For $p\ge q+\tfrac23$, $q>1$, and an absolutely continuous $\g:\R/\Z\to\R$ we have
 $\E(\g)\equiv\infty$ for all $C^{1}$-curves~$\g$.
 To see this, note that we assumed $p<\frac23q+1$ in Theorem~\ref{thm:energy-space} mainly
because neither~\eqref{eq:Wsemi} nor~\eqref{eq:Wsemi'} is not defined for $s\ge1$.
For general $p\ge \frac23q+1$ we nevertheless still have
\begin{equation*}
 \int_{\R/\Z}\int_{-1/2}^{1/2} \frac{\abs{\dg(u+w)-\dg(u)}^q}{\abs{w}^{3p-2q-1}}\d w\d u
 \le C\br{\E(\g)+\E(\g)^{\beta}}.
\end{equation*}
 Applying Brezis~\cite[Prop.~2]{brezis}, the function $\dg$ is constant,
 hence $\g$ lies on a straight line. Therefore, $\g$ cannot be a closed
 $C^{1}$-curve.
\end{remark}


\begin{remark}[Strange energies for $p\in[q+\tfrac23,\tfrac23q+1)$]\label{rem:strange}
 On $p\in[q+\tfrac23,\tfrac23q+1)$, $p,q>0$, see the hatched area in Figure~\ref{fig:range},
 we find the strange behavior that there are no
 closed finite-energy $C^{3}$-curves
 while self-intersections, and in particular corners, are not penalized.
 So piecewise linear curves (polygonals) have finite energy.
 
 The latter can be seen by adapting the calculation
 in Remark~\ref{rem:non-rep}.
 For the former we recall that a closed
 arc-length parametrized $C^2$-curve must have
 positive curvature $\abs{\g''}$ at some point $u_{0}$ and by continuity
 there are $c,\delta>0$ with
 $\abs{\g''}\ge c>0$ on $[u_{0}-\delta,u_{0}+\delta]$.
 As $\g''\perp\dg$ we obtain $\abs{\g''\wedge\dg}=\abs{\g''}\ge c$.
 So $\E(\g)$ is bounded below by
 \begin{align*}
  &\int\limits_{u_{0}-\delta}^{u_{0}+\delta}
  \int\limits_{-\delta}^{\delta}
  \int\limits_{\frac13\abs v}^{\frac23\abs v}
  \abs v^{-3p+2q}\abs{\frac{\T v0}v\wedge\frac{\T 0w}w}^{q} \d w\d v\d u \\
  &=\int\limits_{u_{0}-\delta}^{u_{0}+\delta}
  \int\limits_{-\delta}^{\delta}
  \int\limits_{\frac13\abs v}^{\frac23\abs v}
  \abs v^{-3p+2q}\Bigg|\tfrac{v-w}2\g''(u)\wedge\dg(u)
  +\tfrac{v^{2}}2\int_{0}^{1}(1-\vth_{1})^{2}\g'''(u+\vth_{1}v)\d\vth_{1}\wedge\br{\dg(u)+\tfrac w2\g''(u)} \\
  &\qquad{}-\tfrac{w^{2}}2\int_{0}^{1}(1-\vth_{2})^{2}\g'''(u+\vth_{2}w)\d\vth_{1}\wedge\br{\dg(u)+\tfrac v2\g''(u)} \\
  &\qquad{}+\tfrac{v^{2}w^{2}}2\iint_{[0,1]^{2}}(1-\vth_{1})^{2}(1-\vth_{2})^{2}\g'''(u+\vth_{1}v)\wedge\g'''(u+\vth_{2}w)\d\vth_{1}\d\vth_{2}
  \Bigg|^{q} \d w\d v\d u \\
  &\ge \delta
  \bv{\tilde c-C\delta \lnorm[\infty]{\g'''}^{q}\br{\lnorm[\infty]{\g'''}+\lnorm[\infty]{\g''}+1}^{q}}
  \int_{-\delta}^{\delta}
  \abs v^{-3p+3q+1}\d v.
%  -C\delta^{2}\int\limits_{-\delta}^{\delta}\abs v^{3q-3p+2}\d v
 \end{align*}
 Lessening $\delta>0$, the square bracket is positive.
 This gives $\E(\g)=\infty$.
\end{remark}


\section{\secB}

The arguments here are quite similar as
for the tangent-point energies~\cite{blatt-reiter3},
however, we provide full proofs for the readers' convenience.

Using Theorem~\ref{thm:energy-space} together with
the Arzel\`a-Ascoli theorem, we see that sets of 
curves in $\Cia(\mathbb R / \mathbb Z, \mathbb R^n)$
with a uniform bound on the energy are sequentially compact in $C^1$.
%The next proposition will help us to show that the limits we get
%are injective curves:
To this end we need the following result.

\begin{proposition}[Uniform bi-Lipschitz estimate]\label{prop:bi-Lipschitz}
 For every $M < \infty$ and~\eqref{eq:sub-critical} there is a constant $C(M,p,q)>0$ such 
 that every curve $\gamma\in \Cia(\mathbb R / \mathbb Z, \mathbb R^n)$
 parametrized by arc-length with
 \begin{equation}\label{eq:TP<M}
  \E(\gamma) \leq M
 \end{equation}
 satisfies the bi-Lipschitz estimate
 \begin{equation}\label{eq:bil}
  \abs{u-v} \leq C(M,p,q) \abs{\gamma(u) - \gamma(v)}
  \qquad \text{for all } u,v \in \mathbb R / \mathbb Z.
 \end{equation}
\end{proposition}

The proof is based on the following lemma.
%We will give an easy proof that essential boils down to combining 
%the regularity we get form Proposition~\ref{prop:conttang} with
%a subtle scaling argument. The following lemma will be one of the
%essential parts in the proof.
To be able to state it,
we set for two arc-length parametrized curves $\gamma_i : I_i \rightarrow \mathbb R$, $i=1,2$,
$I_1,I_2$ open intervals,
\begin{align*}
 \E(\gamma_1, \gamma_2) &:= \E(\gamma_1) + \E(\gamma_2) + {} \\
% &{}\qquad +
% \int_{I_1} \int_{I_2} \br{\frac{\dist(\ell_1(t), \gamma_2(s))^q }{|\gamma_1(s)-\gamma_2(t)|^p} 
% + \frac{\dist(\ell_2(t),\gamma_1(s))^q }{|\gamma_2(s)-\gamma_1(t)|^p}}\d s\d t, \\
 &{}\qquad +\iiint\limits_{I_1^2\times I_{2}}
 \frac{\abs{\dg_{1}(u_{1})}\abs{\dg_{1}(u_{2})}\abs{\dg_{2}(u_{3})}}{\RR(\g_{1}(u_{1}),\g_{1}(u_{2}),\g_{2}(u_{3}))}
 {\d u_{1}\d u_{2}\d u_{3}} \\
 &{}\qquad +\iiint\limits_{I_1\times I_{2}^{2}}
 \frac{\abs{\dg_{1}(u_{1})}\abs{\dg_{2}(u_{2})}\abs{\dg_{2}(u_{3})}}{\RR(\g_{1}(u_{1}),\g_{2}(u_{2}),\g_{2}(u_{3}))}
 {\d u_{1}\d u_{2}\d u_{3}}.
\end{align*}

%The proof of Proposition~\ref{prop:bi-Lipschitz} is based on

%We then have
\begin{lemma}\label{lem:last}
 Let $\alpha \in (0,1)$. For $\mu>0$ we let
 $M_\mu$ denote the set of all pairs $(\gamma_1, \gamma_2)$
 of curves $\gamma_i \in \Cia([-1,1],\R^n)$ 
 satisfying
 \begin{enumerate}
  \item $|\gamma_1(0) - \gamma_2 (0)| = 1$,
  \item $\g_{1}'(0)\perp\br{\gamma_1(0)-\gamma_2(0)} \perp \gamma_2'(0)$,
  \item $\|\gamma_{i}'\|_{C^{0,\alpha}} \leq \mu$, \qquad$i=1,2$.
 \end{enumerate}
 Then there is a $c=c(\a,\mu)>0$ such that
 \begin{equation*}
  \E(\gamma_1, \gamma_2) \geq c
  \qquad\text{for all }(\gamma_1, \gamma_2) \in M_\mu.
 \end{equation*}
\end{lemma}

\begin{figure}
 \centering
 \includegraphics[scale=.5]{dist-0.pdf}
 \caption{A pair of curves $(\g_1,\g_{2})\in M_{\mu}$
  defined in Lemma~\ref{lem:last}.
  Note that the arcs $\g_{1}$, $\g_{2}$ cannot intersect each other.}
\end{figure}

\begin{proof}
 It is easy to see that $\E(\gamma_1, \gamma_2)$
 is zero if and only if both $\gamma_1$ and $\gamma_2$ 
 are part of one single straight line.
 We will show that $\E(\cdot,\cdot)$ attains its minimum
 on $M_\mu$. As $M_\mu$ does not contain straight lines by~(i), (ii),
 this minimum is strictly positive which thus proves the lemma.
 
 Let $(\gamma_1^{(n)},\gamma_2^{(n)})$ be a minimizing sequence in 
 $M_\mu$, i.e.\@ we have
 \begin{equation*}
  \lim_{n\rightarrow \infty}\E (\gamma_1^{(n)},\gamma_2^{(n)})
  = \inf_{M_\mu} \E(\cdot,\cdot).
 \end{equation*}
 Subtracting $\g_{1}(0)$ from \emph{both} curves, i.e.\@ setting
 \begin{equation*}
  \tilde \gamma_i^{(n)}(\tau) := \gamma_i^{(n)}(\tau) - \gamma_{1}(0),
  \qquad i = 1,2,
 \end{equation*}
 and using Arzel\`a-Ascoli we can pass to a subsequence such that
 \begin{equation*}
  \tilde \gamma_i^{(n)} \to \tilde \gamma_i\qquad\text{in }C^{1}.
 \end{equation*}
 Furthermore, $(\tilde \gamma_1, \tilde \gamma_2) \in M_\mu$
 since $M_\mu$ is closed under convergence in $C^1$.
 Since, by Fatou's lemma, the functional $\E$ is lower semi-continuous
 with respect to $C^1$ convergence, we obtain
 \begin{equation*}
  \E(\tilde \gamma_1, \tilde \gamma_2) 
  \leq \lim_{n\rightarrow \infty} \E (\tilde \gamma_1^{(n)}, \tilde \gamma_2^{(n)})
  = \lim_{n\rightarrow \infty} \E ( \gamma_1^{(n)}, \gamma_2^{(n)})
  =  \inf_{M_\mu} \E(\cdot,\cdot).
  \end{equation*}
\end{proof}

Let us use this lemma to give the

\begin{proof}[Proposition~\ref{prop:bi-Lipschitz}]
Applying Theorem~\ref{thm:energy-space} to~\eqref{eq:TP<M} we obtain
%\begin{equation*}
 $\|\gamma'\|_{C^{0,\alpha}} \leq C(M)$
%\end{equation*}
 for $\alpha = 3\frac {p-1} q -2\in(0,1-\frac1q)$.
As an immediate consequence there is a $\delta=\delta(\a,M)>0$
such that
\begin{equation}\label{eq:bil2}
 |u-v| \leq 2 |\gamma(u)- \gamma(v)|
\end{equation}
for all $u,v \in \mathbb R / \mathbb Z$ with $|u-v| \leq \delta.$
Let now
\[ S := \inf \sets{\!\rule{0ex}{1em}\abs{\gamma(u)-\gamma(v)}}{u,v \in \R/\Z,\abs{u-v}\geq\delta} \leq \tfrac12. \]
 We will complete the proof by estimating $S$ from below.
 Using the compactness of $\set{u,v \in \R/\Z,\abs{u-v}\geq\delta}$,
 there are $s,t \in \mathbb R / \mathbb Z$ with 
$|s-t| \geq \delta$ and
%\begin{equation*}
 $|\gamma(s)-\gamma(t)| =S$.
%\end{equation*}
If now $|s-t| =\delta$ we obtain
\begin{equation*}
 2 S = 2|\gamma (s) - \gamma(t)| \refeq[\ge]{bil2} \delta
\end{equation*}
and hence
\begin{equation*}
 |u-v| \leq \tfrac12 \le \frac S\delta \le \frac{\abs{\g(u)-\g(v)}}{\delta(\a,M)}
\end{equation*}
for all $u,v \in \mathbb R / \mathbb Z$ with $|u-v| \geq \delta$.
This proves the proposition in this case.
%
If on the other hand $|s-t| > \delta$ then we infer using the
minimality of $|\gamma(s)-\gamma(t)|$
\begin{equation*}
 \gamma'(s) \perp (\gamma(s) - \gamma(t)) \perp \g'(t).
\end{equation*}
We define for $\tau \in [-1,1]$
\begin{equation*}
 \gamma_1 (\tau) := \frac 1 S \gamma(s+ S\tau)
 \qquad\text{and}\qquad
 \gamma_2(\tau) := \frac 1 S \gamma(t + S \tau).
\end{equation*}
Since 
%\begin{equation*}
$\| \gamma'_i\|_{C^{0,\alpha}} \leq \| \gamma'\|_{C^{0,\alpha}}
{\le C(M)}$
%\end{equation*}
we may apply Lemma~\ref{lem:last} which yields
\begin{equation*}
 \E(\gamma_1,\gamma_2) \geq c(\a,M) >0.
\end{equation*}
Together with 
%\begin{equation*}
 $\E (\gamma_1, \gamma_2) \leq S^{3p-2q-3} \E (\gamma)$
%\end{equation*}
this leads to
\begin{equation*}
 S \geq \left( \frac {c(\a,M)}{\E(\gamma)}\right)^{\frac 1 {3p-2q-3}}
 \ge\br{\frac{c(\a,M)}{M}}^{\frac1{3p-2q-3}}.
\end{equation*}
Hence,
$\abs{u-v} \le\tfrac12\le\frac{\abs{\g(u)-\g(v)}}{2S}
   \le C(M,p,q)\abs{\g(u)-\g(v)}$
   for all $u,v \in \mathbb R / \mathbb Z$ with $|u-v| \geq \delta$.
\end{proof}

We are now in the position to prove the
compactness result which is crucial both to
the existence of minimizers
in any knot class
and to the self-avoiding behavior of the energies.

\begin{proposition}[Sequential compactness]\label{prop:sequentiallycompact}
 For each $M< \infty$ the set
 \begin{equation*}
  A_M:=\sets{\gamma \in \Cia(\mathbb R / \mathbb Z, \mathbb R^n)}{ 
  \E(\gamma) \leq M}
 \end{equation*}
 is sequentially compact in $C^1$ up to translations.
\end{proposition}

\begin{proof}
 By Theorem~\ref{thm:energy-space}
 there are $C(M)< \infty$ and $\a=\a(p,q)>0$ such that
% \begin{equation*}
  $\|\dg\|_{C^\alpha} \leq C(M)$
% \end{equation*}
 for all $\g\in A_M$ and hence
 \begin{equation*}
  \|\tilde \g\|_{C^{1,\alpha}} \leq C(M) + 1
 \end{equation*}
 where $\tilde \g (u) := \g(u)-\g(0)$.
% Furthermore, from
% Proposition~\ref{prop:bi-Lipschitz} we get
% the bi-Lipschitz estimate
% \begin{equation*}
%  |u-v| \leq C(M,p,q) |\gamma(u)-\gamma(v)|
% \end{equation*}
% for all $u,v \in \mathbb R / \mathbb Z$.
 By Proposition~\ref{prop:bi-Lipschitz},
 the bi-Lipschitz estimate~\eqref{eq:bil} holds.
 
 Let now $\g_n \in A_M$. Then
 \begin{equation*}
  \|\tilde \g_n\|_{C^{1,\alpha}} \leq C(M) + 1
 \end{equation*}
 and hence after passing to suitable subsequence we have
% \begin{equation*}
  $\tilde \g_n \rightarrow \g_{0}$
% \end{equation*}
 in $C^1$. Since $\g_n$ was parametrized by arc-length,
 $\g_0$ is still parametrized by arc-length and
% still
% \begin{equation*}
%  |u-v| \leq C(M,p,q) |\gamma_{0}(u) - \gamma_{0}(v)|
% \end{equation*}
% for all $u,v \in \mathbb R / \mathbb Z$.
 the bi-Lipschitz estimate carries over to $\g_{0}$.
 So, especially, $\g_{0} \in \Cia(\mathbb R / \mathbb Z, \mathbb R^n)$.
 From lower semi-continuity with respect 
 to $C^1$ convergence we infer
 \begin{equation*}
  \E(\g_0) \leq \liminf_{n\rightarrow \infty} \E(\g_n) \leq M.
 \end{equation*}
 So $\g_0 \in A_M$.
\end{proof}

We may now pass to the

\begin{proof}[Theorem~\ref{thm:existence}]
 Let $\seqn[k]{\g}\in \Cia$ be a minimal sequence for $\E$ in a given knot class $K$,
 i.e.\@ let
 \begin{equation*}
  \lim_{k\rightarrow \infty} \E(\g_{k}) = \inf_{\Cia\cap K} \E.
 \end{equation*}
 After passing to a subsequence and suitable translations, we hence get by Proposition~\ref{prop:sequentiallycompact} 
 a ~$\g_0\in \Cia$ with $\g_k\to\g_0$ in $C^1$.
 As the intersection of every knot class with $C^1$
 is an open set in $C^1$~\cite[Cor.~1.5]{blatt:isot}
 (see~\cite{reiter:isot} for an explicit construction),
 the curve $\g_0$ belongs to the same knot class
 as the elements of the minimal sequence $\seqn[k]{\g}$. 
 The lower semi-continuity of $\E$ furthermore implies
 \begin{equation*}
  \inf_{\Cia\cap K} \E\leq \E(\g_0)
  \leq  \lim_{n\rightarrow \infty} \E(\g_n)=\inf_{\Cia\cap K}\E.
 \end{equation*}
 Hence, $\g_0$ is the minimizer we have been searching for.
\end{proof}

By the same reasoning one derives the existence of a global minimizer of~$\E$.

Let us conclude this section by deriving that the generalized integral Menger curvature
are in fact knot energies (in the sub-critical range).

\begin{proposition}[$\E$ is a strong knot energy {\cite[Cor.~2.3]{SM9}}]\label{prop:strong-knot-energy}%\hfill\\
Let~\eqref{eq:sub-critical} hold.\\[-2em]
 \begin{enumerate}
  \item If $\seqn[k]{\g}$ %\subset\Wir$
  is a sequence of embedded $\W$-curves uniformly converging to a non-injective curve $\g_\infty\in C^{0,1}$ parametrized by arc-length then $\E(\g_k)\to\infty$.
  \item For given $E,L>0$ there are only finitely many knot types having a representative with $\E<E$ and $\text{length}=L$.
 \end{enumerate}
\end{proposition}

\begin{proof}
 The first statement immediately follows from the bi-Lipschitz estimate
 in Proposition~\ref{prop:bi-Lipschitz}, as a sequence with bounded energy 
 would be sequentially compact in $\Cia$ and thus cannot uniformly converge 
 to a non-injective curve.   
 
 To show the second statement, let us assume that it was wrong, i.e.,
 that there are curves $\seqn[k]{\g}$ of length $L$,
 all belonging to different knot classes,
 with energy less than $E$. Of course we can assume that
 $L=1$.
 After applying suitable transformations and passing to a subsequence,
 Proposition~\ref{prop:sequentiallycompact} guarantees the existence
 of $\g_{0}\in A_{M}$ with $\g_{k}\to\g_{0}$ in $C^1$.
 Again by~\cite{blatt:isot,reiter:isot}
 this implies that almost all $\g_k$ belong to the
 same knot class as $\g_0$, which is a contradiction.
\end{proof}